% Theme: Nocturne  --  best for one formula, explained in depth
% Build:  ./build.sh example-nocturne.tex   (or compile on Overleaf)
\documentclass[12pt]{beamer}
\usetheme{ignocturne}

\ighandle{@quantum_danesh}
\igseries{Fórmula destacada}
\igissue{07}

\begin{document}

\coverslide{Información cuántica}
  {Todo canal es, en secreto, un estado}
  {El isomorfismo de Choi--Jamio\l{}kowski}

\begin{frame}{La definición}{La fórmula}
Toma un par máximamente entrelazado y pasa \key{una mitad} por el canal
$\Phi$. Lo que sale es la matriz de Choi:
\keyeq{J(\Phi) = (\mathrm{id}\otimes\Phi)\big(\ketbra{\Omega}\big)}
con el estado sin normalizar
\[
  \ket{\Omega} = \sum_{i=1}^{d} \ket{i}\otimes\ket{i}.
\]
Una sola matriz de $d^2\times d^2$ guarda todo lo que hace $\Phi$.
\end{frame}

\begin{frame}{Término a término}{Desglose}
Al desarrollar $\ketbra{\Omega}$ aparece una matriz por bloques:
\[
  J(\Phi) = \sum_{i,j}
    \explain<2->{\ket{i}\!\bra{j}}{qué bloque}
    \otimes
    \explain<3->{\Phi\big(\ket{i}\!\bra{j}\big)}{qué le hace $\Phi$ a $\ket{i}\!\bra{j}$}
\]
\uncover<4->{El bloque $(i,j)$ es la salida para la entrada $\ket{i}\!\bra{j}$.
Por linealidad, esas $d^2$ salidas \hl{determinan $\Phi$ por completo}.}
\end{frame}

\begin{frame}{El diccionario}{Propiedades}
Cada propiedad del canal se vuelve una propiedad de una matriz:
\begin{align*}
  \Phi \text{ completamente positivo} &\iff \mhl{J(\Phi) \succeq 0}\\[1mm]
  \Phi \text{ preserva la traza}      &\iff \Tr_{\mathrm{out}} J(\Phi) = \id\\[1mm]
  \Phi \text{ unital}                 &\iff \Tr_{\mathrm{in}} J(\Phi) = \id\\[1mm]
  \text{rango de Kraus de } \Phi      &= \operatorname{rango} J(\Phi)
\end{align*}
\begin{notebox}
La primera línea es la razón de su fama: la positividad completa, una
afirmación sobre todas las extensiones, se reduce a calcular autovalores.
\end{notebox}
\end{frame}

\begin{frame}{Dos extremos}{Ejemplos}
\begin{defbox}[Canal identidad]
$J(\mathrm{id}) = \ketbra{\Omega}$: rango uno y máximamente entrelazado.
Un solo operador de Kraus, $K = \id$.
\end{defbox}
\begin{thmbox}[Totalmente despolarizante]
$\Phi(X) = \Tr(X)\,\id/d$ da $J(\Phi) = \id\otimes\id/d$: rango completo y
nada de entrelazamiento. Necesita $d^2$ operadores de Kraus.
\end{thmbox}
\end{frame}

\begin{frame}{Por qué compensa}{En la práctica}
\begin{itemize}
  \item Ver si un mapa es un canal se reduce a calcular autovalores.
  \item Los autovectores de $J$, reordenados como matrices $d\times d$ y
        escalados por $\sqrt{\lambda_k}$, son operadores de Kraus.
  \item Optimizar sobre canales se vuelve un programa semidefinido en $J$.
\end{itemize}
\[
  J = \sum_k \lambda_k \ketbra{v_k}
  \;\Longrightarrow\;
  K_k = \sqrt{\lambda_k}\;\mhl*{\operatorname{unvec}(v_k)}
\]
\end{frame}

\closingslide{Una matriz, todo el canal}
  {Próxima fórmula: la dilatación de Stinespring. Sígueme para no perderla.}

\end{document}
